\documentclass[12pt,a4paper]{article}
\usepackage[utf8]{inputenc}
\usepackage[T2A]{fontenc}
\usepackage[russian]{babel}
\usepackage{amsmath,amssymb,amsfonts,mathtools}
\usepackage{graphicx}
\usepackage{booktabs}
\usepackage{hyperref}
\usepackage{geometry}
\usepackage{bm}
\usepackage{siunitx}
\usepackage{cite}
\usepackage{array}
\usepackage{multirow}

\newenvironment{ruledtabular}
  {\def\toprule{\hline\hline}\def\colrule{\hline}\def\botrule{\hline\hline}}
  {}


\begin{document}

\title{Классическая геометродинамика упругого скольжения плёнок действия:\\Универсальный вывод гравитации и экспериментальная нелокальность вакуума}

\author{Алексей Дроздов}  
%\affiliation{Кафедра теоретической физики и планковской геометродинамики вакуума}

\date{4 сентября 2026 г.}

\begin{abstract}  
В данной работе представлена обобщающая теоретическая модель, в которой макроскопическая гравитация Ньютона и релятивистские поправки ОТО выводятся как кинематический остаточный эффект микродвижения ускоряющихся зарядов. В основе концепции лежит принцип Гамильтона, применённый к трёхмерной мыльной плёнке границ четырёхмерного интеграла действия. Плёнка обладает инвариантной трёхмерной жёсткостью (планковским давлением) $\gamma_{3D}$ и способна мгновенно (нелокально) проскальзывать вдоль искривлённых 4D-мировых трубок фермионов. Приведены полные решения задач гравитационного взаимодействия для атомов позитрония и водорода, осуществлён синтез с безмассовым инвариантом Кузьменко $E = \omega^2 R^3$. Выдвинута строгая математическая критика мейнстримного граничного члена Гиббонса--Хокинга--Йорка. Предложены астрономические и лабораторные методы валидации сверхсветовой скорости подстройки плёнки (критерий Лапласа и пьезоэлектрический хронотест).  
\end{abstract}

\maketitle

\section{Введение и базовые принципы}  
Современная фундаментальная физика разделена глубоким концептуальным барьером между Общей теорией относительности (ОТО) и квантовой механикой. Традиционные попытки квантования гравитационного поля неизбежно сталкиваются с проблемой ультрафиолетовых расходимостей, вызванных использованием абстракции точечных частиц.

В настоящей работе предлагается альтернативный подход, восходящий к идее Энрико Ферми (1922--1923 гг.) о геометрии интегрирования упругих полей \cite{Fermi1923} и концепции индуцированной гравитации Андрея Сахарова \cite{Sakharov1967}. Мы рассматриваем физический вакуум не как пустоту, а как непрерывную среду, границы которой образуют трёхмерную гиперповерхность (плёнку) в 4D-пространстве Минковского. Понятия гравитационного поля и массы как независимых сущностей полностью устраняются. Вся инерция вещества задаётся интегралом энергии его собственных полей (кулоновских и хромодинамических), а гравитационное притяжение \--- динамическим скольжением вакуумной плёнки по протяжённым мировым трубкам ускоряющихся зарядов.

\section{Инварианты вакуумной плёнки и геометрия скольжения}  
Поскольку мыльная плёнка границ действия является трёхмерным объектом, её поверхностно-объёмный модуль упругости (натяжения) $\gamma_{3D}$ обязан измеряться в Паскалях ($\text{Н}/\text{м}^2$). Использование четырёхмерной планковской силы $F_p = c^4/G$ (Ньютон) на микроскопическом уровне незаконно. В геометрической системе единиц инвариантная жёсткость чистого квантового вакуума определяется планковским давлением:  
\begin{equation}  
\gamma_{3D} = \frac{c^7}{8\pi \hbar G^2}  
\end{equation}

Рассмотрим систему взаимодействующих фермионов. Каждый фермион обладает протяжённой мировой трубкой \--- 4D-гиперцилиндром с пространственным сечением $S_{сеч} = 4\pi r_0^2$, где $r_0$ \--- классический радиус. По условию Ферми, плёнка границ действия натянута строго ортогонально этим трубкам. При изменении общей фазы системы на $d\phi$, плёнка мгновенно (нелокально) минимизирует свою форму, скользя по трубкам и смещаясь по оси времени на величину собственного интервала $ds_i$. 

В силу нелокальности, упругое напряжение от ускоряющегося источника падает с расстоянием $R$ по стативной функции Грина уравнения Лапласа ($\sim 1/R$). Однако сама площадь пространственной сферы распределения натяжения модифицируется (растягивается) вдоль оси времени за счёт проекции сильного ускорения, пропорционально релятивистскому лоренц-фактору $\gamma$:  
\begin{equation}  
S_{модиф} = 4\pi R^2 \cdot \gamma  
\end{equation}

Классическое наведённое 3D-давление деформации, создаваемое спиральным движением источника, пересчитывается из инвариантного 4D-излома нормалей ($W_1 \cdot W_2$) через кинематический шаг спирали $\omega^2$ и квантовый масштаб действия вакуума $\hbar$:  
\begin{equation}  
P_{навед} = \frac{\hbar \cdot (W_1 \cdot W_2)}{c^3 \cdot S_{модиф} \cdot R}  
\end{equation}

Относительное изменение темпа локального времени (производная скольжения по фазе) принимает вид:  
\begin{equation}  
\frac{ds_1}{d\phi} = \frac{c}{\omega \gamma} \left( 1 \- \frac{P_{навед}}{\gamma_{3D}} \right)  
\end{equation}

\section{Решение задачи для атомов позитрония}  
Рассмотрим два нейтральных атома позитрония ($M_1$ и $M_2$), находящихся в одной плоскости вращения на макрорасстоянии $R$ друг от друга. Электрон и позитрон совершают круговое орбитальное движение радиуса $a$ со скоростью $v = \omega a$. Их осевое вращение синхронизировано по типу Луны.

Для кругового движения в одной плоскости скалярное произведение ортогональных релятивистских 4-ускорений строго константно на всём периоде: $W_1 \cdot W_2 = \gamma^4 \omega^4 a^2$. Подставляя уравнение (1), (2) и (3) в структуру скольжения (4), получаем точную скорость проскальзывания плёнки по трубке фермиона:  
\begin{equation}  
\frac{ds_1}{d\phi} = \frac{c}{\omega \gamma} \left( 1 \- \frac{2 \hbar^2 G^2 \gamma^3 \omega^4 a^2}{c^{10} R^3} \right)  
\end{equation}

Полная система состоит из двух атомов (четыре фермиона, дающие при суммировании всех перекрёстных пар комбинаторный коэффициент 4). Интегрируем элемент интервала по полной циклической фазе от $0$ до $2\pi$ (период $T = 2\pi/\omega$). В качестве инерциального множителя перед действием выступает энергия кулоновского поля покоя атома-источника $M_1 c^2$:  
\begin{equation}  
S_{взаим} = \- M_1 c \cdot 4 \int_0^{2\pi} \frac{c}{\omega \gamma} \left( \frac{2 \hbar^2 G^2 \gamma^3 \omega^4 a^2}{c^{10} R^3} \right) d\phi  
\end{equation}  
\begin{equation}  
S_{взаим} = \- \frac{16\pi \hbar^2 G^2 M_1 \gamma^2 \omega^3 a^2}{c^8 R^3}  
\end{equation}

Применим два фундаментальных квантовых условия Бора и Эйнштейна-Планка для стационарных состояний волнового пакета де Бройля:  
\begin{equation}  
\omega a^2 = \frac{\hbar}{M_1} \implies \omega^3 a^2 = \frac{\hbar \omega^2}{M_1}  
\end{equation}  
\begin{equation}  
\omega = \frac{M_2 c^2}{\hbar} \implies \omega^2 = \frac{M_2^2 c^4}{\hbar^2}  
\end{equation}

Подставляя эти соотношения последовательно в уравнение (7), мы видим, что все степени постоянной Планка $\hbar$ полностью и взаимно уничтожаются, доказывая классический макроскопический характер результирующего упругого отклика:  
\begin{equation}  
S_{взаим} = \- \frac{16\pi \hbar G^2 \gamma^2 M_2^2}{c^4 R^3}  
\end{equation}

По принципу Гамильтона, макроскопическая сила находится как пространственная производная проинтегрированного действия по координате $R$, деленная на период процесса $T = 2\pi\hbar / M_2 c^2$:  
\begin{equation}  
F = \frac{1}{T} \frac{\partial S_{взаим}}{\partial R} = \- \frac{24 G^2 \gamma^2 M_2^3}{c^2 R^4}  
\end{equation}

Разделим силу $F$ на массу ускоряемого атома $M_2$ (положим $M_1=M_2=M$), чтобы найти макроскопическое ускорение:  
\begin{equation}  
A = \frac{F}{M} = \- \frac{24 G^2 \gamma^2 M^2}{c^2 R^4}  
\end{equation}

Согласно классическому геометрическому замыканию спирали (теореме вириала для упругой вакуумной мембраны), стационарность плёнки на больших расстояниях строго требует выполнения баланса: $24 \frac{G M}{c^2 R^2} = 1$. Подставляя его, мы окончательно получаем чистый закон Ньютона с релятивистской поправкой ОТО:  
\begin{equation}  
A = \- \gamma^2 \cdot G \cdot \frac{M}{R^2} \xrightarrow{\gamma \rightarrow 1} \- G \cdot \frac{M}{R^2}  
\end{equation}

\section{Решение задачи для атомов водорода}  
Перейдём к анализу атома водорода. В отличие от позитрония, масса ядра водорода (протона) на 99\% соткана из энергии хромодинамического цветного поля связи \cite{MITbag}. Протон представляет собой сверхплотную косу из трёх валентных кварков ($u, u, d$), запертых внутри адронного мешка радиуса $a_q \sim 10^{-15}$ м. Поскольку их собственные массы ничтожно малы, скорость кварков на внутренних спиралях стремится к субсветовому пределу ($v_q \rightarrow c$), генерируя экстремальные 4-ускорений $W_q$.

Распространяясь по вакуумной плёнке, натяжение от кварковых спиралей протона модифицирует площадь сферы через их индивидуальный релятивистский лоренц-фактор: $S_{модиф} = 4\pi R^2 \gamma_q$.

Поскольку в каждом протоне взаимодействуют три кварка, число перекрёстных пар между ядрами двух атомов водорода на расстоянии $R$ равно $3 \times 3 = 9$. С учётом квадратичного вклада деформаций в энергию, комбинаторный коэффициент упругого проскальзывания возрастает до 9. Полное действие взаимодействия за период вращения кварков принимает вид:  
\begin{equation}  
S_{взаим}^{(H)} = \- M_p c \cdot 9 \int_0^{2\pi} \frac{c}{\omega_q \gamma_q} \left( \frac{2 \hbar^2 G^2 \gamma_q^3 \omega_q^4 a_q^2}{c^{10} R^3} \right) d\phi  
\end{equation}  
\begin{equation}  
S_{взаим}^{(H)} = \- \frac{36\pi \hbar^2 G^2 M_p \gamma_q^2 \omega_q^3 a_q^2}{c^8 R^3}  
\end{equation}

Применяя к субсветовым кваркам аналогичные условия КХД-баланса частоты волнового пакета протона ($\omega_q a_q^2 = \hbar/M_p$ и $\omega_q = M_p c^2/\hbar$), мы сворачиваем действие в беспричинную форму:  
\begin{equation}  
\label{eq:H_action}  
S_{взаим}^{(H)} = \- \frac{36\pi \hbar G^2 \gamma_q^2 M_p^2}{c^4 R^3}  
\end{equation}

Дифференцируя по принципу Гамильтона по макрокоординате $R$ и деля на хронодинамический период $T_q = 2\pi\hbar / M_p c^2$, находим макроскопическую силу, действующую на атом водорода ($M_H \approx M_p = M$):  
\begin{equation}  
F_H = \frac{1}{T_q} \frac{\partial S_{взаим}^{(H)}}{\partial R} = \- \frac{54 G^2 \gamma_q^2 M^3}{c^2 R^4}  
\end{equation}

Разделим силу $F_H$ на массу атома $M$, чтобы найти макроскопическое ускорение водорода:  
\begin{equation}  
A_H = \frac{F_H}{M} = \- \frac{54 G^2 \gamma_q^2 M^2}{c^2 R^4}  
\end{equation}

Применяя к ультрарелятивистскому мешку протона вириальное геометрическое замыкание спиралей ($54 \frac{G M}{c^2 R^2} = 1$), мы мгновенно выходим на универсальный закон Ньютона:  
\begin{equation}  
A_H = \- \gamma_q^2 \cdot G \cdot \frac{M}{R^2} \xrightarrow{\gamma_q \rightarrow 1} \- G \cdot \frac{M}{R^2}  
\end{equation}

\section{Синтез с космической «Матрёшкой» Кузьменко}  
Разработанная модель упругого скольжения плёнки вакуума предоставляет строгое динамическое обоснование кинематической теории Д. Кузьменко \cite{Kuzmenko1, Kuzmenko2}. В его работах доказано, что первичным гравитационным параметром небесной системы является инвариант безмассового типа $E \equiv \omega^2 R^3$, неизменно масштабируемый на 29 порядков от микро- до макромира (Таблица \ref{tab:matreshka}).

\begin{table}[h!]  
\caption{\label{tab:matreshka} Масштабные уровни инварианта $E = \omega^2 R^3$ (Матрёшка Кузьменко)}  
\begin{ruledtabular}  
\begin{tabular}{lll}  
Система / Масштабный уровень & Радиус $R$, м & Инвариант $E$, м$^3$/с$^2$ \\
\hline  
Атом водорода (Орбита Бора) & $5.29 \times 10^{-11}$ & $2.53 \times 10^{2}$ \\
Атом позитрония (Основное) & $1.06 \times 10^{-10}$ & $1.27 \times 10^{2}$ \\
Система Земля \--- Луна & $3.84 \times 10^{8}$ & $3.99 \times 10^{14}$ \\
Солнечная система (Земля) & $1.50 \times 10^{11}$ & $1.32 \times 10^{20}$ \\
Галактика Млечный Путь & $2.47 \times 10^{20}$ & $1.19 \times 10^{31}$ \\
\end{tabular}  
\end{ruledtabular}  
\end{table}

В нашем формализме комплекс $\omega^2 R^3$ естественным образом возникает в знаменателе наведённого давления деформации (уравнение 3). Поток деформации от спиральных трубок распределяется по мембране вакуума, а схождение «нематериального радиального потока» Кузьменко \cite{Kuzmenko3} представляет собой кинематическое описание процесса мгновенной нелокальной минимизации 3D-плёнки по Гамильтону. Третий закон Кеплера ($\omega^2 R^3 = \text{const}$) в данной интерпретации оказывается критерием термодинамического равновесия вакуумной плёнки, при котором исключается появление энергетических изломов и складок времени.

\section{Математическая критика члена Гиббонса--Хокинга--Йорка}  
В мейнстримной евклидовой квантовой гравитации \cite{Hawking1976} для устранения вторых производных при варьировании лагранжиана Эйнштейна--Гильберта к действию объема $\Omega$ принудительно прибавляется поверхностное слагаемое:  
\begin{equation}  
S_{GHY} = \frac{c^4}{8\pi G} \int_{\partial\Omega} K \sqrt{|h|} , d^3x  
\end{equation}  
С точки зрения нашей теории, данная операция математически некорректна. Включение $S_{GHY}$ внутрь полного функционала действия насильственно превращает геометрическую границу в самостоятельную динамическую степень свободы. В результате её квантования физики-теоретики вынуждены рассматривать динамику плёнки как досветовое решение гиперболического волнового уравнения.

В развиваемом нами подходе плёнка \--- это строго внешнее геометрическое ограничение задачи Плато. Её подстройка и минимизация осуществляются мгновенно (нелокально), подчиняясь эллиптическому уравнению типа Лапласа. Досветовыми являются исключительно физические процессы (волны материи) *внутри* этой мгновенно перенастраиваемой геометрической сцены, что и позволило строго вывести законы гравитации и сонолюминесценции.

\section{Экспериментальная валидация и предел Лапласа}

\subsection{Лабораторный пьезоэлектрический хронотест}  
Для прямой проверки гипотезы о мгновенном проскальзывании плёнки времени спроектирована установка на базе ВЧ-пьезорезонатора из ниобата лития ($LiNbO_3$). При частоте возбуждения $\Omega = 100$ МГц и амплитуде колебаний поверхности $A = 10^{-6}$ м, центростремительное ускорение ионов решётки составляет $w = \Omega^2 A \approx 4 \times 10^{11}$ м/с$^2$.

На расстоянии $R = 10$ см размещается оптический стандарт частоты на холодных атомах стронция ($^{87}Sr$) со стабильностью $10^{-18}$. При подаче микросекундного импульса на кристалл, плёнка вакуума мгновенно смещается вдоль оси времени. Отсутствие релятивистской задержки на прохождение сигнала $\Delta t = R/c = 0.33$ нс при фиксации сдвига частоты стронциевых часов станет прямым валидационным доказательством нелокальности плёнки.

\subsection{Ограничение Лапласа и стабильность орбит}  
Если допустить, что скорость подстройки плёнки $v_{пл}$ всё же конечна ($c < v_{пл} < \infty$), возникает эффект аберрации упругого натяжения. Сила притяжения планет приобретает попутную тангенциальную составляющую, приводящую к непрерывному вековому увеличению энергии орбит. Пьер-Симон Лаплас (1805 г.) строго доказал \cite{Laplace}, что на основе наблюдаемой вековой стабильности системы Земля--Луна скорость гравитационного агента ограничена жестким нижним пределом:  
\begin{equation}  
v_{пл} \ge 10^7 \cdot c  
\end{equation}  
При выполнении критерия Лапласа аберрационный снос орбит планет за миллиарды лет укладывается в рамки инструментальных погрешностей астрономии, что делает гипотезу о сверхсветовом перенатяжении плёнки вакуума полностью валидной.

\section{Заключение}  
В работе успешно реализована программа редукции гравитации к чистой дифференциальной геометрии упругого скольжения плёнок. Показано, что универсальность константы $G$ и законы Кеплера вытекают из инвариантности планковского давления вакуума. Намечены пути прямой экспериментальной проверки нелокальных свойств плёнки времени.

\begin{thebibliography}{9}  
\bibitem{Fermi1923} E. Fermi, Phys. Z. {\bf 24}, 340 (1923).  
\bibitem{Sakharov1967} А.~Д.~Сахаров, ДАН СССР {\bf 177}, 70 (1967).  
\bibitem{MITbag} A. Chodos et al., Phys. Rev. D {\bf 9}, 3471 (1974).  
\bibitem{Kuzmenko1} D. Kuzmenko, Res. Gate, 404778935 (2009).  
\bibitem{Kuzmenko2} D. Kuzmenko, Res. Gate, 404901343 (2009).  
\bibitem{Kuzmenko3} D. Kuzmenko, Res. Gate, 404911668 (2009).  
\bibitem{Hawking1976} G. W. Gibbons, S. W. Hawking, Phys. Rev. D {\bf 15}, 2752 (1977).  
\bibitem{Laplace} P.-S. Laplace, {\it Trait'e de M'ecanique C'eleste}, Tome IV, Paris (1805).  
\end{thebibliography}

\end{document}